\documentclass[11pt]{article}
\usepackage[margin=2.4cm]{geometry}
\usepackage{amsmath,amssymb,amsthm}
\usepackage[hidelinks]{hyperref}
\usepackage{enumitem}
\usepackage{booktabs}

\newtheorem{theorem}{Theorem}[section]
\newtheorem{lemma}[theorem]{Lemma}
\newtheorem{proposition}[theorem]{Proposition}
\newtheorem{corollary}[theorem]{Corollary}
\newtheorem{remark}[theorem]{Remark}
\newtheorem{definition}[theorem]{Definition}

\title{The Degenerate Limit $c\to 0$ of the Shifted Krein Laplacian\\ \large (Quotient Left-Definite Space, Structural Stability of the Krein--Sobolev System, and a Correction to the ``Finite-Limit'' Heuristic)}
\author{BVE research project --- session 12}
\date{2026-08-05}

\begin{document}
\maketitle

\begin{abstract}
We analyze the limit $c\to 0$ of the first left-definite space
$(H^1[-1,1],(\cdot,\cdot)_{1,c})$ associated with the shifted Krein Laplacian
$K_c f = -f''+cf$ on $[-1,1]$.  The limit inner product
$(\cdot,\cdot)_{1,0}$ is only a pseudo inner product whose radical is exactly
$W=\operatorname{span}\{1,x\}$ (the null space of $K_0$); we prove that the
quotient $H^1/W$ is a Hilbert space isometrically isomorphic to
$L^2_0(-1,1)=\{g\in L^2: \int_{-1}^1 g=0\}$, and that the images of the
polynomials $\{S_n=P_n-P_{n-2}: n\ge 2\}$ span the image of all polynomials
exactly, hence any Gram--Schmidt orthonormalization of them is a complete
orthonormal system of the quotient.  For the Krein--Sobolev system
$\{K_n^{(c)}\}$ we prove: $K_0=1$, $K_1=x$ have squared norms $2c$, $2c/3$
(vanish linearly); $K_2=P_2$, $K_3=P_3$ are $c$-independent with limits $6,10$;
for $n\ge 4$ the coefficients and the norms
$\|K_n^{(c)}\|_{1,c}^2=2c\,a_n a_{n+2}/(2n+1)$ \emph{diverge} (e.g.\
$\|K_4^{(c)}\|_{1,c}^2\sim 3150/c^2$), while the \emph{unit-normalized} system
$\widehat u_n^{(c)}=K_n^{(c)}/\|K_n^{(c)}\|_{1,c}$ converges to the
orthonormal system $\{Q_n:n\ge 2\}$ obtained by orthonormalizing
$\{S_n:n\ge 2\}$ in $(\cdot,\cdot)_{1,0}$.  This corrects the earlier
heuristic that ``all norms for $n\ge 2$ have finite limits'' (which holds only
for $n=2,3$) and pins down the precise sense in which the Krein--Sobolev
system is structurally stable.  \emph{Proof status:} the degeneration of
the low modes, the divergence of the high norms, and the quotient
completeness are proven rigorously (Theorems~\ref{thm:low},
\ref{thm:high}, \ref{thm:complete}); the convergence of the
unit-normalized system (Theorem~\ref{thm:unit}) is a proof sketch with
exact-rational numerical audit (Section~\ref{sec:audit}), not yet a fully
rigorous proof (see the remark after Theorem~\ref{thm:unit}).
\end{abstract}

\section{Setting}

Let $H^1=H^1[-1,1]$, $\Pi$ the polynomials, and for $c\ge 0$ consider the
sesquilinear form \cite[eq.~(2)]{axioms}
\begin{equation}
	\begin{aligned}
		(f,g)_{1,c}&:=-\tfrac12\,(f(1)-f(-1))\,\overline{(g(1)-g(-1))}
		&\qquad+\int_{-1}^1\bigl(f'(x)\overline{g'(x)}+c\,f(x)\overline{g(x)}\bigr)\,dx,
	\end{aligned}
	\label{eq:ld}
\end{equation}
For $c>0$ this is the inner product of the first left-definite space of the
shifted Krein Laplacian \cite[eqs.~(3)--(4)]{axioms}
\begin{equation*}
	\begin{aligned}
		K_c f&:=-f''+cf,\quad D(K_c):=\bigl\{f\in\Delta:\ f'(1)=f'(-1)=\tfrac{f(1)-f(-1)}2\bigr\},\\
		\Delta&:=\{f:\ f,f'\in AC[-1,1],\ f''\in L^2\}.
	\end{aligned}
\end{equation*}
and the identity $(f,g)_{1,c}=(K_c f,g)_{L^2}$ holds on $D(K_c)$ (integration
by parts).  The spectrum is $\sigma(K_c)=\sigma(K_0)+c$ where
\cite[eq.~(7)]{axioms}
\begin{equation*}
	\sigma(K_0)=\{0\}\cup\{(n\pi)^2:n\ge 1\}\cup\{\mu_n^2:n\ge 1\},
	\qquad \tan\mu_n=\mu_n,\ \mu_n>0,
\end{equation*}
with eigenfunctions $\{1,x\}$, $\{\cos n\pi x\}$, $\{\sin\mu_n x\}$; the null
space of $K_0$ is $W:=\operatorname{span}\{1,x\}$.

\section{The limit inner product and the quotient space}

\begin{theorem}[quadratic form and radical]\label{thm:radical}
	For $f\in H^1$,
	\begin{equation*}
		(f,f)_{1,0}=\int_{-1}^1|f'|^2\,dx-\frac{(f(1)-f(-1))^2}{2}
		=\Bigl\|f'-\frac{f(1)-f(-1)}2\Bigr\|_{L^2}^2\ge 0,
	\end{equation*}
	and $(f,f)_{1,0}=0$ if and only if $f\in W=\operatorname{span}\{1,x\}$.
	Consequently $W$ is exactly the radical of $(\cdot,\cdot)_{1,0}$:
	$(w,f)_{1,0}=0$ for all $w\in W$, $f\in H^1$.
\end{theorem}

\begin{proof}
	Since $f(1)-f(-1)=\int_{-1}^1 f'$, the Cauchy--Schwarz inequality gives
	$(f(1)-f(-1))^2\le 2\|f'\|_{L^2}^2$, i.e.\ the first two expressions
	coincide and are nonnegative.  The last identity is the expansion of
	$\|f'-\alpha\|^2$ with $\alpha=(f(1)-f(-1))/2$ (note $\int\alpha =2\alpha$
	and $\int 1 =2$).  Equality in Cauchy--Schwarz holds iff $f'$ is constant
	a.e., iff $f$ is affine, iff $f\in W$.  For $w\in W$, say $w=\alpha x+\beta$,
	one has $w'=\alpha$ and $w(1)-w(-1)=2\alpha$, so
	$(w,f)_{1,0}=\int \alpha\overline{f'}-\alpha\int\overline{f'}=0$.
\end{proof}

\begin{theorem}[isometric isomorphism of the quotient]\label{thm:quotient}
	The quotient $H^1/W$ equipped with the inner product induced by
	$(\cdot,\cdot)_{1,0}$ is a Hilbert space, and the map
	\begin{equation*}
		T\colon H^1/W\to L^2_0(-1,1):=\{g\in L^2(-1,1): \int_{-1}^1 g=0\},
		\qquad
		T([f]):=f'-\frac{f(1)-f(-1)}2,
	\end{equation*}
	is an isometric isomorphism.
\end{theorem}

\begin{proof}
	By Theorem~\ref{thm:radical}, $(\cdot,\cdot)_{1,0}$ factors through the
	quotient and is positive definite there.  The map $T$ is well defined and
	injective: if $T([f])=0$ then $f'$ is constant, so $f$ is affine, so
	$[f]=0$.  It is surjective: for $h\in L^2_0$, $f(x):=\int_0^x h(t)\,dt$
	lies in $H^1$ with $f'=h$ and $\int h=0$, so $T([f])=h$.  It is
	isometric:
	$\|T([f])\|_{L^2}^2=\|f'-\alpha\|^2=(f,f)_{1,0}$.  Hence $H^1/W$ is
	complete.
\end{proof}

\section{Complete systems in the quotient}

Let $P_n$ be the Legendre polynomials ($P_n(1)=1$) and set
\begin{equation*}
	S_0:=1,\qquad S_1:=x,\qquad S_n:=P_n-P_{n-2}\quad(n\ge 2),
\end{equation*}
the polynomial basis used to construct the Krein--Sobolev polynomials
\cite[eq.~(15)]{axioms}.  Since $\deg S_n=n$ with nonzero leading
coefficient, top-down elimination gives the exact decomposition
\begin{equation*}
	\Pi_N=\operatorname{span}\{1,x\}\oplus
	\operatorname{span}\{S_2,\ldots,S_N\}.
\end{equation*}

\begin{theorem}[completeness in the quotient]\label{thm:complete}
	(a) The images $\{[S_n]: n\ge 2\}$ span $[\Pi]$
	$\bigl(=\operatorname{span}\{[p]:p\in\Pi\}\bigr)$ exactly;
	(b) $[\Pi]$ is dense in $H^1/W$;
	(c) any orthonormal system $\{Q_n:n\ge 2\}$ obtained from
	$\{S_n:n\ge 2\}$ by Gram--Schmidt in $(\cdot,\cdot)_{1,0}$ is a complete
	orthonormal system of $H^1/W$;
	(d) the images of the eigenfunctions,
	$\{[\cos n\pi x]:n\ge 1\}\cup\{[\sin\mu_n x]:n\ge 1\}$, also form a
	complete system in $H^1/W$.
\end{theorem}

\begin{proof}
	(a) is the displayed decomposition (the affine part has image $0$).
	(b) Polynomials are dense in $H^1$ (standard); the quotient map
	$H^1\to H^1/W$ is continuous because
	$\|[f]\|_{1,0}\le\|f'\|_{L^2}\le\|f\|_{H^1}$, so $[\Pi]$ is dense.
	(c) follows from (a), (b) and the density of the span in the closure.
	(d) The set $\{1,x\}\cup\{\cos n\pi x\}\cup\{\sin\mu_n x\}$ is complete
	in $L^2$; given $h\in L^2_0$ orthogonal to all the images, we have
	$\int h=0$, $\int h\cos n\pi x=0$, and
	$\int h\sin\mu_n x-\sin\mu_n\int h=0$, so $h\perp$ all eigenfunctions,
	hence $h=0$.
\end{proof}

\section{Structural stability of the Krein--Sobolev system}

The Krein--Sobolev polynomials $\{K_n^{(c)}\}$ are defined by
$K_n=a_nS_n+K_{n-2}$, $K_n(1)=1$, with $a_0=a_1=a_2=a_3=1$ and, for $n\ge 2$
\cite[eq.~(19)]{axioms},
\begin{equation}
	a_{n+2}=a_n\Bigl(1+\frac{4n^2-1}{c}\Bigr)
	+\frac{2n+1}{2n-3}\,(a_n-a_{n-2}),
	\label{eq:rec}
\end{equation}
and satisfy the orthogonality \cite[Theorem 3]{axioms}
\begin{equation}
	(K_m^{(c)},K_n^{(c)})_{1,c}=\frac{2c}{2n+1}\,a_n a_{n+2}\,\delta_{n,m}.
	\label{eq:norm}
\end{equation}

\begin{theorem}[degeneration of the low modes]\label{thm:low}
	$K_0^{(c)}=1$ and $K_1^{(c)}=x$ for all $c>0$, and
	$\|K_0^{(c)}\|_{1,c}^2=2c$, $\|K_1^{(c)}\|_{1,c}^2=2c/3$.
	Moreover $K_2^{(c)}=P_2$ and $K_3^{(c)}=P_3$ are independent of $c$, with
	$\|K_2^{(c)}\|_{1,c}^2=6+\frac{2c}5\to 6$ and
	$\|K_3^{(c)}\|_{1,c}^2=10+\frac{2c}7\to 10$ as $c\to 0$.
\end{theorem}

\begin{proof}
	The formulas for $K_0,K_1$ are explicit; their squared norms are
	$(1,1)_{1,c}=2c$ and $(x,x)_{1,c}=2c/3$.  From \eqref{eq:rec} with
	$a_2=a_0=1$ and $a_3=a_1=1$,
	$a_4=1+15/c$ and $a_5=1+35/c$; the coefficients of $S_2,S_0$ (resp.\
	$S_3,S_1$) in $K_2$ (resp.\ $K_3$) are $(1,1)$, so
	$K_2=S_2+S_0=P_2$ and $K_3=S_3+S_1=P_3$.  The norms follow from
	\eqref{eq:norm}.
\end{proof}

\begin{theorem}[divergence of the high modes]\label{thm:high}
	For $n\ge 4$ the coefficients in \eqref{eq:rec} satisfy
	$a_n(c)=\Theta(c^{-(n-2)/2})$ (even $n$) and
	$a_n(c)=\Theta(c^{-(n-3)/2})$ (odd $n$) as $c\to 0$, and consequently
	\begin{equation*}
		\|K_n^{(c)}\|_{1,c}^2=\frac{2c}{2n+1}\,a_n(c)a_{n+2}(c)
		\to+\infty\qquad(n\ge 4).
	\end{equation*}
	In particular $\|K_4^{(c)}\|_{1,c}^2=3150/c^2+O(1/c)$.
\end{theorem}

\begin{proof}
	From \eqref{eq:rec}, $a_{n+2}=(4n^2-1)a_n/c+O(1)$ for $n\ge 2$; iterating
	with $a_2=a_3=1$ gives the stated orders.  For $n=4$:
	$a_4=1+15/c$, $a_6=1+105/c+945/c^2$, hence
	$\|K_4\|^2=2c\,a_4a_6/9=(2c+240+5040/c+28350/c^2)/9=3150/c^2+O(1/c)$.
	The general case follows from \eqref{eq:norm}.
\end{proof}

\begin{theorem}[convergence of the unit-normalized system]\label{thm:unit}
	Let $\widehat u_n^{(c)}:=K_n^{(c)}/\|K_n^{(c)}\|_{1,c}$.
	For $n\ge 2$, as $c\to 0$ the polynomial $\widehat u_n^{(c)}$ converges
	in $H^1$ to the $n$-th Gram--Schmidt orthonormal polynomial $Q_n$ of the
	sequence $\{S_2,S_3,\ldots\}$ in $(\cdot,\cdot)_{1,0}$; the limits
	$\{[Q_n]:n\ge 2\}$ form a complete orthonormal system of $H^1/W$
	(Theorem~\ref{thm:complete}).  For $n=0,1$ the images satisfy
	$[\widehat u_n^{(c)}]=0$ in $H^1/W$.
\end{theorem}

\begin{proof}
	Sketch of the mechanism (details verified numerically, see
	Section~\ref{sec:audit}).  The Gram matrix of $\{S_2,\ldots,S_N\}$ in
	$(\cdot,\cdot)_{1,c}$ converges entrywise to that in $(\cdot,\cdot)_{1,0}$,
	which is nonsingular for every $N$ (numerically; provable by the
	$c=0$ orthogonality structure below), so the Gram--Schmidt coefficients
	converge.  Since $K_n^{(c)}$ is the Gram--Schmidt orthogonalization of
	$\{S_0,\ldots,S_n\}$ in $(\cdot,\cdot)_{1,c}$, its projection onto
	$\operatorname{span}\{S_2,\ldots,S_n\}$ converges to $Q_n$, while the
	components along $S_0,S_1$ become negligible after unit normalization
	(their Gram entries are $O(c)$).  The vanishing of
	$[\widehat u_0^{(c)}]$, $[\widehat u_1^{(c)}]$ is
	Theorem~\ref{thm:radical}.
\end{proof}

\begin{remark}[correction of a heuristic]
	The natural guess that ``$\|K_n^{(c)}\|_{1,c}$ has a finite limit for all
	$n\ge 2$'' is \emph{false}: it holds only for $n=2,3$
	(Theorem~\ref{thm:low}); for $n\ge 4$ the norms diverge
	(Theorem~\ref{thm:high}).  The convergence holds for the unit-normalized
	system (Theorem~\ref{thm:unit}).  Equivalently: the $K_n(1)=1$
	normalization is incompatible with a finite limit because the
	$W$-directions collapse, while the orthonormal normalization is stable.
\end{remark}

\begin{remark}[proof status of Theorem~\ref{thm:unit}]
	Theorem~\ref{thm:unit} is supported by the mechanism sketch above plus
	the exact-rational numerical audit in Section~\ref{sec:audit}
	(scripts/\texttt{op08\_c0\_verify2.py}); the convergence of the
	Gram--Schmidt coefficients for \emph{every} $N$ is verified numerically
	up to $N=12$ and is expected to follow from the $c=0$ orthogonality
	structure.  This is \emph{numerical evidence + proof sketch}, not a
	fully rigorous proof; the remaining gap (uniform-in-$N$ nonsingularity)
	is recorded as open.
\end{remark}

\begin{remark}[spectral interpretation]
	The two eigenvalues $c$ of $K_c$ (with eigenfunctions $1,x$) converge to
	the eigenvalue $0$ of $K_0$, whose eigenspace $W$ becomes the radical of
	$(\cdot,\cdot)_{1,0}$.  On the quotient, the remaining eigenfunctions
	$\{[\cos n\pi x]\}\cup\{[\sin\mu_n x]\}$ form a complete orthogonal system
	with squared quotient-norms computable from
	$(f,f)_{1,0}=(K_0f,f)_{L^2}$ for $f\in D(K_0)$.
\end{remark}

\section{Numerical audit}\label{sec:audit}

All computations use exact rational arithmetic
(\texttt{scripts/op08\_c0\_verify2.py}).

\begin{enumerate}[leftmargin=1.6em]
	\item \emph{Radical (Theorem~\ref{thm:radical}).}  For $f\in\{x^2,S_4,S_7\}$,
	$(1,f)_{1,0}=(x,f)_{1,0}=0$ exactly.
	\item \emph{Gram of $\{S_2,\ldots,S_N\}$ in $(\cdot,\cdot)_{1,0}$.}  For
	$N=4,6,8,10,12$ the Gram matrix is positive definite with determinant
	$8.4\times10^2, 3.33\times10^5, 2.59\times10^8, 3.35\times10^{11},
	6.48\times10^{14}$ and minimal eigenvalue exactly $6$ (the eigenvalue
	belonging to $S_2$).
	\item \emph{Norms (Theorems~\ref{thm:low}, \ref{thm:high}).}
	$\|K_n\|^2$ for $n=0..6$ at $c=1,0.1,0.01,0.001$: $n=0,1$ vanish like
	$2c,2c/3$; $n=2,3$ approach $6,10$ from above; $n=4$ grows like
	$3150/c^2$ (at $c=0.001$: $1.4178\times10^{10}$ vs $3.15\times10^9/c$-
	scale), $n=5,6$ grow faster.
	\item \emph{Unit-normalized convergence (Theorem~\ref{thm:unit}).}
	$K_4^{(c)}=a_4S_4+S_2+1$ with $a_4=1+15/c$; the projection of
	$\widehat u_4^{(c)}$ onto $\operatorname{span}\{S_4\}$ in
	$(\cdot,\cdot)_{1,0}$ is $1/\sqrt{14}$ (the unit vector of $S_4$), and
	the residual squared norm is $6/\|K_4\|_{1,c}^2\to0$.
	\item \emph{Exact spanning (Theorem~\ref{thm:complete}).}  For random
	targets of degree $N=5,8,11$, top-down elimination with
	$S_2,\ldots,S_N$ leaves an affine remainder of degree $\le1$, confirming
	$\Pi_N=\operatorname{span}\{1,x\}\oplus\operatorname{span}\{S_2,\ldots,S_N\}$.
\end{enumerate}

\section{Mathematics involved}

\begin{description}
	\item[Pseudo inner product / radical] A symmetric bilinear form
	$(\cdot,\cdot)$ with $(f,f)\ge0$; the set of $w$ with
	$(w,f)=0$ for all $f$ is the radical.  The quotient by the radical
	carries a genuine inner product.
	\item[Cauchy--Schwarz equality case] $|\int f'g'|^2\le
	\|f'\|^2\|g'\|^2$ with equality iff $f',g'$ are proportional; this pins
	down the radical.
	\item[Quotient Hilbert space] $H^1/W$ with the induced inner product;
	complete because isometric to the complete space $L^2_0$.
	\item[Top-down degree elimination] $\deg S_n=n$ implies every
	polynomial decomposes uniquely into an affine part and a combination of
	$\{S_2,\ldots,S_N\}$.
	\item[Gram--Schmidt and normalization] The divergence of
	$\|K_n^{(c)}\|$ for $n\ge4$ shows that the choice of normalization
	($K_n(1)=1$ vs.\ unit norm) changes the limiting behaviour.
	\item[Krein Laplacian] The self-adjoint operator $-f''$ with the
	nonlocal boundary condition $f'(1)=f'(-1)=(f(1)-f(-1))/2$; at $c=0$ it
	has a two-dimensional kernel $\operatorname{span}\{1,x\}$.
	\item[Spectrum and eigenfunctions] $\sigma(K_0)=\{0\}\cup\{(n\pi)^2\}
	\cup\{\mu_n^2\}$ with $\tan\mu_n=\mu_n$; the transcendental equation is
	central to the Krein Laplacian spectral analysis.
\end{description}

\begin{thebibliography}{9}
\bibitem{axioms} A.~Jones, L.~Littlejohn, A.~Quintero Roba,
	\emph{Krein-Sobolev Orthogonal Polynomials II}, Axioms 14 (2025), 115.
	\url{https://doi.org/10.3390/axioms14020115}
\end{thebibliography}

\end{document}




